\section{Further results on persistence and the long run}
\label{app:moved-ak}

This section collects results on persistence and the long run that the main text does
not need: the order-selection block, the transition-strength corollary and a pointer to the
two-sided statement; and, from the sections on the general results, post-entry non-termination and the selection
boundary, the economic mechanism behind the basins, persistence from primitives, the
pathwise ordering under basin preservation and the comparative statics of lock-in. Each
statement is the one the formalization map cites.

\subsection{The order-selection block (\Cref{sec:mechanism})}

\begin{proposition}[Order-selected steady state]
\label{prop:order-selection}
Let $\basin$ be preserved by the equilibrium closed loop $G_E(x)=\F(x,R_E(x))$,
and suppose that the closed loop is non-decreasing on $\basin$ (a higher state
today never leads to a lower state tomorrow, which holds whenever the
equilibrium sanction rule $R_E$ is non-decreasing), and that $R_E$ is
\emph{continuous} on $\basin$. Fix a start $x_0\in\basin$.
\begin{enumerate}[label=(\roman*)]
\item If $x_0\le R_E(x_0)$, the political state path converges to the
  \emph{least} steady state of $G_E$ at or above $x_0$.
\item If $R_E(x_0)\le x_0$, it converges to the \emph{greatest} steady state of
  $G_E$ at or below $x_0$.
\end{enumerate}
In either case the realized sanction converges to $R_E$ evaluated at that limit.
The two cases are exhaustive, and nothing is assumed about how many steady states
$\basin$ contains.
\end{proposition}

\begin{corollary}[Coincidence of the two limits]
\label{cor:order-selection-identified}
If in addition the transition satisfies strict directional adjustment
\textup{(F5)}, the realized sanction and the political state converge to the same
value.
\end{corollary}

Continuity of $R_E$ is the requirement; uniqueness of the maximizer of
$J_S$ is one sufficient condition for it
(\Cref{app:moved-ad}).

Order selection need not imply attraction of a whole basin. The selected limit may
vary with the initial condition. For a monotone continuous closed loop on an
invariant interval, every path converges to one common limit if and only if the
interval contains a unique steady state; equivalently, $G_E$ lies strictly off the
diagonal on either side of that state. The next result obtains uniqueness by
deriving the interval from the first steady state rather than imposing an
additional condition.

\begin{proposition}[Attraction on a derived interval]
\label{prop:derived-interval}
Under the hypotheses of \Cref{prop:order-selection} (an invariant basin, a
non-decreasing closed loop and a continuous sanction rule, with no level
condition), let $x_s$ be
the \emph{least} steady state of $G_E$ at or above the basin pivot $\bar x$. Then
every equilibrium path started anywhere in $[\bar x,x_s]$ converges to $x_s$:
a single limit, common to every start in that interval. Under strict directional
adjustment \textup{(F5)} the realized sanction converges to $x_s$ as well.
\end{proposition}

In plain terms: initial conditions determine the direction of movement and select
an extremal steady state on that side. Different initial states in the same
invariant interval need not converge to a common point: unless the interval
runs only as far as the first steady state, in which case they do.

\paragraph{What the three results say about sanctions regimes.}
A steady state of $G_E$ is a political state at which the sender's optimal
sanction reproduces the entrenchment it inherits: escalation has no further
return, and a lower sanction would concede entrenchment the sender wants kept. The regime
sustains itself without either government changing its behavior.

\Cref{prop:order-selection} says that the inherited political state selects,
at date zero and within a single equilibrium, which self-sustaining regime is
reached (path dependence, not indeterminacy), and
\Cref{prop:derived-interval} that up to the first such regime above the pivot
every start reaches it. None of the three results delivers a rate: they say
where a sanctions regime ends up, not how long it takes
(\Cref{app:moved-ah}).

\subsection{The transition-strength corollary (\Cref{sec:compstat})}

\begin{corollary}[Transition strength and shock duration]
\label{cor:shock-transition-strength}
Consider two transition laws $\F_1$ and $\F_2$ satisfying
\textup{(F1)--(F4)} and evaluated at the same threshold $0\le\bar x$, below
which the first rises under full sanctions, $x<\F_1(x,1)$ for
$x\in[0,\bar x]$ (so $\bar x<1$). If
\[
\F_1(x,1)\le\F_2(x,1)
\qquad\text{for every }x\in[0,\bar x],
\]
then
\[
\eta_{\F_1}(\bar x)\le\eta_{\F_2}(\bar x),
\qquad
\widehat N_{\F_2}(\bar x)
\le\widehat N_{\F_1}(\bar x).
\]
If the transition inequality is strict throughout $[0,\bar x]$, the
drift inequality is strict. The duration inequality remains weak
because distinct drift bounds can have the same integer ceiling.

Under the affine law,
\[
\eta_\delta(\bar x)=\delta(1-\bar x),
\qquad
\widehat N_\delta(\bar x)
=\left\lceil
\frac{\bar x}{\delta(1-\bar x)}
\right\rceil.
\]
Hence $\delta_1\le\delta_2$ implies
$\widehat N_{\delta_2}(\bar x)\le
\widehat N_{\delta_1}(\bar x)$.
\end{corollary}

\subsection{The same-equilibrium two-sided statement (\Cref{sec:relief-exit})}

The statement, \Cref{prop:two-sided-one-equilibrium}, is in
\Cref{app:dual-basin-economy}, with the definition of the dual-basin economy.

\subsection{Post-entry non-termination and the selection boundary}
\label{app:nontermination}

Fix a date $T$ with $x_T^E\in\basin$ and restart the stationary
equilibrium path from $x_T^E$. The proof of
\Cref{thm:persistence} applies to this tail and gives
\[
s_{T+n}^E\ge\bar x,
\qquad
s_{T+n}^E\le s_{T+n+1}^E
\quad\text{for every }n\ge0.
\]
For a cutoff $q<\bar x$, define
\[
\tau_q^E(T)
:=\inf\{n\ge0:s_{T+n}^E\le q\},
\qquad \inf\varnothing:=+\infty.
\]
The strict cutoff inequality and the inherited floor make the set
empty, so $\tau_q^E(T)=+\infty$. A weak cutoff is insufficient: a path
that stays exactly at its floor reaches that weak threshold. Only the floor is
used, so the conclusion needs only the level condition, the sanction rule at or
above the state on the basin (\path{equilibriumTop_tailF_no_relief_of_ratchet});
neither A1 nor the target condition enters.

Condition (E4) supplies the optimization content. At every tail date
and against every feasible sanction $r\in[0,1]$,
\[
\begin{split}
&U_S\bigl(x_{T+n}^E,\alpha(x_{T+n}^E),s_{T+n}^E\bigr)
 +\beta_S V_S\!\left(F(x_{T+n}^E,s_{T+n}^E)\right)\\
&\qquad\ge
 U_S\bigl(x_{T+n}^E,\alpha(x_{T+n}^E),r\bigr)
 +\beta_S V_S\!\left(F(x_{T+n}^E,r)\right).
\end{split}
\]
Thus non-termination is selected from an action set that includes
relief and every reduction of the sanction. Without state-local dynamic-argmax
uniqueness, this conclusion is specific to the greatest selection: the greatest
maximizer stays above the floor, while another tied maximizer need not.
With uniqueness, every maximizer equals the realized action and every
other sanction is strictly worse. Lean verifies the tail
restart, floor, empty first-passage set, optimality comparison, and
uniqueness upgrade in
\path{Sanctions/Theory/LockIn/PostEntryNonTermination.lean}.

\subsection{Economic mechanism}
\label{sec:mechanism}
\label{sec:mechanism:forces}

Political lock-in combines two forward-looking incentives. A sender
that values a more entrenched successor state weighs current sanctions
against the entrenchment they leave. For feasible sanctions $s_L<s_H$, the dynamic sender
prefers $s_H$ to $s_L$ exactly when
\[
\beta_S\!\left[
\Vs(\F(x,s_H))-\Vs(\F(x,s_L))
\right]
>
\Us(x,a,s_L)-\Us(x,a,s_H).
\]
The discounted continuation gain must exceed the current payoff loss.
This is the direct effect of the sender's foresight; it underlies the
comparison of dynamic and myopic sanctions at a fixed target policy in
\Cref{thm:dominance}.

The target faces the corresponding intertemporal trade-off. More
compliance can reduce the current sanction but also erode a political
constituency valued by the incumbent. A Bellman-optimal target may
therefore comply less than a myopic target. The two incentives point the same way:
lower compliance raises the sender's sanction (the indirect effect of the target's lower compliance), and the sender's
continuation value raises the sanction relative to its myopic best response.
Whether they \emph{amplify} one another is a separate question, and no result
below signs an interaction primitive (\Cref{app:moved-ad}).

\paragraph{Formation of the high basin.}
\label{sec:mechanism:basin}
Under \Cref{ass:flip}, the target and sender flip sets are upper
intervals, and \Cref{prop:tipping} defines their infima as political
pivots. Above their joint pivot, the sender's incentive supports a
sanction floor and the target's current-plus-future objective favors
the high-sanction direction. The target condition supplies preference
evidence for the latter claim. Basin invariance itself follows from the
sender's equilibrium floor and transition monotonicity.

\subsection{Primitive applications and scope}
\label{sec:existence:nonvacuity}

The conditions of \Cref{thm:persistence} are attained in an economy where
continuation values affect observable choices (\Cref{thm:b2-benchmark},
\Cref{app:witness:benchmark}); its description and its two robustness
extensions are in \Cref{app:moved-ad}.

\subsubsection{Persistence from primitives}
\label{sec:existence:primitive-persistence}

The high-basin theorem (\Cref{thm:persistence}) starts from an equilibrium and value sets satisfying its conditions.
A narrower class, whose economies we call the corner economies, permits both to be derived from payoff and transition
primitives. Its transition law satisfies the model's maintained
continuity, monotonicity, diagonal, and strict-adjustment conditions,
together with a lower bound on its sanction slope and a two-sided bound on the
corner section $x\mapsto\F(x,1)$. No upper bound on the sanction slope and no bound on the
curvature of $\F$ are imposed, and
it contains non-affine realized recursions, including members whose
curvature in the sanction exceeds half the lower bound on their sanction slope.

\begin{proposition}[Primitive persistence over a general transition]
\label{prop:primitive-persistence}
For every corner economy (an economy of the general-$\F$ class above) there exists a
stationary feedback Stackelberg equilibrium, lying in invariant sets constructed from the
economy's own coefficients, such that from every initial state in the
high basin, the political state remains in the basin at every date, the
realized sanction never falls below the basin pivot, and the realized
sanction sequence is non-decreasing.
\end{proposition}

No equilibrium, value function, policy, single-crossing condition,
sanction floor, or path property is assumed in this proposition. Those
objects are constructed from inequalities among the primitives. The
gain in primitive content comes with a precise limitation: the
equilibrium sanction is maximal at every state. The result establishes
persistence, but its non-decreasing sanction path is flat rather than a
state-dependent ratchet.

No curvature bound is needed because the sender's choice is settled by
monotonicity rather than by concavity. Every sender value in the constructed
invariant set rises with the state at least at a fixed rate, every unit of sanction
raises the next state by at least the lower bound on the transition's sanction slope, and this class's
implacability condition makes the discounted value of that rise exceed the largest
current cost of the unit. The sender's forward-looking objective is therefore
strictly increasing in the sanction at every state and at every compliance level
the target can choose, and full sanction is its unique maximizer, however
curved the transition is. The constructed invariant set for the sender accordingly restricts only
the level and the slope of the sender's value. If the corner section is in
addition semiconcave with some modulus $M\ge0$, the equilibrium sender value is
semiconcave as well, with a modulus proportional to $M$; no curvature of $\F$ in
the sanction is needed for this. This class contains a member whose transition is
semiconcave with modulus $1/5$, and with no smaller modulus, both in the
sanction and along the corner section; that modulus exceeds half the lower bound on the member's
sanction slope, $3/20$,
and the proposition holds there (\Cref{app:formalization}).

The off-corner boundary of this result is described in \Cref{app:moved-ad}
and \Cref{sec:discussion:scope}.

\subsection{Pathwise ordering under basin preservation}
\label{app:persistence:pathwise}

The corollary below compares the myopic equilibrium $M$ with the equilibrium that a
local invariant-set certificate of the high-basin route constructs. It is an
existence statement about that constructed equilibrium, not a statement about
every equilibrium satisfying the hypotheses of \Cref{thm:persistence}.

\begin{corollary}[Pathwise ordering under basin preservation]
\label{cor:pathwise-dominance}
Let the transition follow the general law \eqref{eq:transition}, let
$\bar x_T,\bar x_S\in[0,1]$ with $\bar x<1$, and take the following
local invariant-set certificate: invariant sets $\mathcal C_T\times\mathcal C_S$ of value functions
on which the sender's dynamic argmax and the target's argmax are unique, as in
\Cref{thm:existence}, and on which the joint Bellman operator satisfies the row
bounds of \Cref{thm:existence} with non-negative coefficients and
$\beta_T+L_{TS}<1$, $\beta_{S,\mathrm{row}}+L_{ST}<1$; \textup{(P2)}
and \textup{(P3)} of \Cref{thm:persistence}; and a sanction floor
$\bar s\in[0,1]$, a slope $m_S$ and a rate $d$ such that, for every value
pair and every $x\in[\bar x_S,1]$, the sender's stage payoff at the induced
target policy rises at rate at least $m_S$ on $[0,\max\{x,\bar s\}]$, the
transition's sanction slope is at most $d$ there, and $\beta_SK_Sd<m_S$. This
rate form implies the pivot-star slope condition \textup{(P5)}; the target condition
\textup{(P4)} is not needed. Suppose in addition that the sanction rule of the
myopic equilibrium satisfies, at every $x\in\basin$, the ceiling
$\sigM(x,\alphaM(x))\le\max\{x,\bar s\}$, and that it preserves the basin:
\begin{equation}
\sigM(x,\alphaM(x))\in\basin
\qquad\text{for every }x\in\basin.
\label{eq:myopic-basin-preservation}
\end{equation}
Then these hypotheses construct an equilibrium $E=(\Vt,\Vs,\alpha)\in\mathcal E$
with $\Vt\in\mathcal C_T$ and $\Vs\in\mathcal C_S$ such that, at every
$x\in\basin$,
\[
\sigM(x,\alphaM(x))\le\sigDynTop_{\Vs}(x,\alpha(x)),
\]
and from every common initial state $x_0\in\basin$ the
realized sanctions satisfy $s^M_t\le s^E_t$ for every $t\ge0$.
\end{corollary}

The corollary is stated at the scope of the certificate theorem that proves
it: the statewise inequality compares the myopic ceiling with the sender's
floor and uses neither \Cref{thm:dominance} nor the compliance ordering
$\alpha\le\alphaM$, the certificate reads no condition on how the target ranks
sanction levels, and the extra hypothesis \eqref{eq:myopic-basin-preservation}
is substantive because two different state paths are compared
(\Cref{app:moved-ah}; the proof is \Cref{app:proof-pathwise}). Both
paths are evolved by the general transition.

\paragraph{Proof of \Cref{cor:pathwise-dominance}.}
\label{app:proof-pathwise}
The corollary is the certificate theorem \path{pathwise_dominance_generalF};
its hypotheses are the fields of the core certificate
\path{StandardComplianceLocalCoreCertificate}, listed in the statement, and basin
preservation \eqref{eq:myopic-basin-preservation}, a separate hypothesis of the
theorem. Write
$\hat x=\max\{x,\bar s\}$.

\emph{Step 1: the equilibrium, its floor and its direction.} Row contraction
gives $E=(\Vt,\Vs,\alpha)\in\mathcal E$ with $\Vt\in\mathcal C_T$,
$\Vs\in\mathcal C_S$ and $\alpha=\alpha_{\Vt,\Vs}$, which is (P1) by construction
(the formal statement exports membership of the invariant sets)
(\path{exists_feedbackStackelbergTop_from_valueClass2_with_membership}). At
$x\in\basin$ the rate form gives the transition-adjusted slope condition
\eqref{eq:F-adjusted-sender-chord} on $[0,\hat x]$
(\path{sender_transition_adjusted_chord_of_upperChordOn_lowerSlope}). Together
with the $K_S$-Lipschitz bound (P2) it makes $J_S(x,\alpha(x),\cdot\,;\Vs)$
strictly increasing on $[0,\hat x]$
(\path{senderDynObjective_strictMono_of_transition_adjusted_chord_and_V_lipschitz}),
as in \Cref{lem:single-crossing}(b) with $\bar s$ in place of $\bar x$. So the dynamic best response
is at least $\hat x$: $R_E(x)\ge\hat x\ge x$
(\path{SenderPersistenceBasinTop.level_dyn_top}), and since $R_E\le1$, $R_E$ maps
$\basin$ into itself. By (P3), $R_E$ is non-decreasing on $\basin$
(\path{dynamicTopRealizedRule_mono_of_IDOFamilyOn}). Both facts are packaged in
\path{exists_equilibrium_senderPersistenceBasinTop_generalF}.

\emph{Step 2: statewise.} At $x\in\basin$ the myopic ceiling and Step~1 give
$\sigM(x,\alphaM(x))\le\hat x\le R_E(x)=\sigDynTop_{\Vs}(x,\alpha(x))$
(\path{ruleDom_of_dynamic_floor_and_myopic_ceiling}).

\emph{Step 3: pathwise.} Let $R_M(x)=\sigM(x,\alphaM(x))$, and let $x^M_t$ and
$x^E_t$ be the paths from $x_0\in\basin$ under $R_M$ and $R_E$, both evolved by
$\F$. Each rule maps $\basin$ into itself: $R_M$ by
\eqref{eq:myopic-basin-preservation} and $R_E$ by Step~1. If $y\in\basin$ and
$R(y)\ge\bar x$, then $\F(y,R(y))\ge\F(\bar x,\bar x)=\bar x$ by (F2)--(F4),
so both paths stay in $\basin$. Suppose $x^M_t\le x^E_t$. Step~2 at $x^M_t$,
monotonicity of $R_E$ on $\basin$, (F3) and (F2) give
\[
x^M_{t+1}=\F\bigl(x^M_t,R_M(x^M_t)\bigr)\le\F\bigl(x^M_t,R_E(x^E_t)\bigr)
\le\F\bigl(x^E_t,R_E(x^E_t)\bigr)=x^E_{t+1}.
\]
So $x^M_t\le x^E_t$ at every date
(\path{rulePathF_le_of_rule_le_and_upper_monoOn_basin}), and
$s^M_t=R_M(x^M_t)\le R_E(x^M_t)\le R_E(x^E_t)=s^E_t$
(\path{ruleSanctionsF_le_of_rule_le_and_upper_monoOn_basin}).

The statewise inequality compares the myopic ceiling with the sender's floor;
it does not go through \Cref{thm:dominance} or the compliance ordering. Basin
preservation is used only in Step~3, where the two paths are at different
states. \qed

\subsection{Comparative statics of lock-in}
\label{sec:compstat:lockin}

Consider two configurations with the same transition law and initial
state. If the second sanction rule is non-decreasing and
pointwise no lower than the first on the high basin, transition
monotonicity implies
\[
x_t^{(1)}\le x_t^{(2)}\qquad\text{for every }t.
\]
When both paths satisfy the conditions of \Cref{thm:persistence}, their limiting
lock-in states have the same order. Thus any primitive or equilibrium
change that raises the sanction rule raises the state path, subject
to the stated monotonicity and common-transition conditions.

Adjustment speed and patience enter differently. With the sanction rule held
fixed, the steady-state set does not depend on a positive affine adjustment
rate; and at a fixed continuation value and target policy greater sender
patience weakly raises the dynamic best response, the corner economies permitting a
total comparison (\Cref{prop:equilibrium-patience-feedback},
\Cref{tab:compstat-windows}, \Cref{app:moved-ah}).
